\section{Finite-horizon physical realization}
\label{sec:finite-realization}

We now reconnect the scalar theorem to physical states.  This is the first point at which an actual Haxell packing is chosen.

Let $\eta>0$ be the resource margin from Proposition~\ref{prop:wide-start}.  By \eqref{eq:role-summable}, choose a rank threshold $B_{\mathrm{res}}$ such that
\[
\sum_{Q>B_{\mathrm{res}}}\operatorname{roleMass}(Q)<\eta.
\tag{7.1}\label{eq:resource-tail}
\]
Also choose a threshold $B_{\mathrm{ns}}$ beyond which every WideStart survivor row has cardinality at least the quota denominator $L_P$ of Proposition~\ref{prop:haxell-quota}.  This exists because the original row has size $\gg Q/\log Q$ and Proposition~\ref{prop:wide-start} preserves a fixed positive fraction.

Apply Theorem~\ref{thm:scalar-terminality} with
\[
B_\dagger=\max\{B_{\mathrm{res}},B_{\mathrm{ns}}\}.
\tag{7.2}\label{eq:global-start-cutoff}
\]
Thus the WideStart base $B$ and the entire scalar first-hit construction are chosen only after the resource tail and the non-small packing onset have been fixed.

Fix one scalar-legal horizon $X\ge B$.  There are only finitely many current rows between $B$ and $X$.  For each such current $Q$, delete from the original row $A_Q$ every atom conflicting with any label of the fixed WideStart.  Let $\operatorname{Surv}(Q)$ denote the resulting row.  Proposition~\ref{prop:wide-start} gives
\[
|\operatorname{Surv}(Q)|
\ge \sigma |A_Q|
\ge \kappa\sigma\frac Q{\log Q}
\tag{7.3}\label{eq:survivor-density}
\]
for all rows in the late regime.

Put an edge between two surviving atoms whenever they cannot coexist in one physical state.  Proposition~\ref{prop:signed-inverse-capabilities} gives a uniform maximum degree bound.  Apply Proposition~\ref{prop:haxell-quota} once to the finite union of rows through $X$.  We obtain a globally independent atom pool $P_X$ with
\[
|P_X\cap\operatorname{Surv}(Q)|
\ge \rho_P|\operatorname{Surv}(Q)|
>\rho\frac Q{\log Q}
\tag{7.4}\label{eq:packed-density}
\]
by the choice \eqref{eq:rho-choice}.  Therefore the actual row projection
\[
P_X(Q)=P_X\cap\operatorname{Surv}(Q)
\]
belongs to the already-defined universal family $\operatorname{Adm}_\rho(Q)$ from \eqref{eq:adm-universal}.  No scalar constant is chosen after seeing $P_X$.

\subsection{Local physical lifts}

The proof of Theorem~\ref{thm:uniform-profile} used full finite solution relations: distinct-value occurrence lifts in the diversity branch and full finite equal-value subset families in the redundancy branch.  Apply those same relations to the actual row $P_X(Q)$.  Since every element is an actual signed-inverse atom, each local target rectangle is now represented by literal finite collections of physical atoms.

Because $P_X$ is globally independent in the interval-conflict graph, all atoms in the pool are pairwise disjoint and nonadjacent.  On a path, every finite subset of such a pool is therefore a physical state.  In addition, every surviving atom was deleted against the complete WideStart obstacle, so it is compatible with \emph{every} WideStart label, not merely with one selected branch.

Consequently the ordinary composability pullback for the entire finite chain maps to $\operatorname{Phys}^n$ through the genuine $n$-ary multiplication domain $D_n$.  The strictification principle of Section~\ref{sec:states-filtration} then identifies the observable composition with the literal physical composition on this finite diagram.  Induction through the current events up to $X$ produces a finite physical response
\[
R_X:\mathbf1\longrightarrow
\operatorname{Tor}_{E_X}(\gamma_X)\times I_X
\tag{7.5}\label{eq:finite-response}
\]
whose projection to the displayed target is surjective.  Here $\gamma_X$ is the transported image of the fixed WideStart centre $\gamma_B$.

\subsection{Resource control}

Let $S_0$ be any WideStart label.  Proposition~\ref{prop:wide-start} gives
\[
W(S_0)<2-\eta.
\]
Every current response at rank $Q$ uses reciprocal value at most $\operatorname{roleMass}(Q)$.  By the subadditive filtration and \eqref{eq:resource-tail}, every edge label $S$ of \eqref{eq:finite-response} satisfies
\[
W(S)
\le W(S_0)+\sum_{B<Q\le X}\operatorname{roleMass}(Q)
<2.
\tag{7.6}\label{eq:physical-W-less-2}
\]
This is the last consumer of the conflict graph, the actual Haxell packing, the $n$-ary compatibility witness, and the resource series.

\subsection{The realization surjection}

Let $\operatorname{Horizon}_B$ be the set of scalar-legal finite cutoffs after the fixed WideStart base.  For $X\in\operatorname{Horizon}_B$, let $\operatorname{Lift}(X)$ be the finite-data solution set consisting of a quota packing through $X$, the local occurrence-lift witnesses, the $n$-ary physical factorization, and the resulting response \eqref{eq:finite-response} satisfying \eqref{eq:physical-W-less-2}.  The preceding argument proves
\[
\operatorname{Lift}(X)\ne\varnothing
\qquad(X\in\operatorname{Horizon}_B).
\]
Define the dependent sum
\[
\operatorname{Real}
=
\sum_{X\in\operatorname{Horizon}_B}\operatorname{Lift}(X).
\]
Then the projection is surjective:
\begin{equation}\label{eq:real-surjection}
\boxed{
\operatorname{Real}\twoheadrightarrow\operatorname{Horizon}_B.}
\end{equation}
Different horizons may use unrelated finite packings.  No coherent infinite choice, inverse limit, or compactness argument occurs.
